\usepackage[normalem]{ulem}
\normalem

\makeatletter
\newif\ifreviewchanges
\@ifundefined{reviewchanges}{\reviewchangesfalse}{\reviewchangestrue}
\makeatother

\newcommand{\reviewbubble}[1]{%
  \tikz[baseline=(mark.base)]{%
    \node[
      circle,
      draw=black,
      fill=white,
      inner sep=0.3pt,
      minimum size=1.35em,
      font=\scriptsize
    ] (mark) {#1};%
  }%
}
\newcommand{\reviewfootnote}[2]{\footnote{\raggedright\textbf{Comment #1.} #2\par}}
\newcommand{\qid}[1]{%
  \ifreviewchanges
    \textcolor{blue!70!black}{\uline{\textbf{[#1]}}}\ %
  \else
    \textbf{[#1]}\ %
  \fi
}
\newcommand{\pdfAnnComment}[3]{%
  \textbf{Original annotated PDF #1.}
  \textbf{Exact annotation:} \emph{#2}
  Response: #3%
}

\newcommand{\annCommentOne}{\textbf{Exact available note:}
\emph{Co presne jsou zaporne tlaky? "negative liquid pressure means postiive
suction relative to the gas/reference pressure".}  Context: clarify the pressure
reference convention, while preserving the Richards sign convention used by OGS.}
\newcommand{\annCommentTwo}{\textbf{Exact available note:}
\emph{Jake jsou typicke chyby/jejich odhady jednotlivych experimentu a vyslednych
velicin?}  Context: record which uncertainty estimates are active policies and
which remain gates.}
\newcommand{\annCommentThree}{\textbf{Exact available note:}
\emph{Lehce offotpic, jak funguje modelovani 2d rezu pro elasticitu, ktera ma
materialove parametry dane z 3d modelu?}  Context: explain the 2D slice
interpretation of fixed 3D-derived mechanical parameters.}
\newcommand{\annCommentFour}{\textbf{Exact available note:}
\emph{je mozne uplne vyhodit zavislost na teplote z OGS? lepsi by asi bylo na to
nesahat, ale pokud by nam umeli dat jen HM model, tak by to mohlo neco malicko
usetrit.}  Context: keep the checked TRM XML form but make the active isothermal
reading explicit.}
\newcommand{\annCommentFive}{\textbf{Exact available note:}
\emph{Chceme nejak integrovat data z "crackmeteru" a podobnych? (spise ne?)}
Context: keep these streams as validation/prior gates unless numeric residual
exports are supplied.}
\newcommand{\annCommentSix}{\textbf{Exact available note:}
\emph{Co znamena S\_lr pro ruzna mereni? S\_e z van Genutchtena, ale permeabilita
je skrze S\_l. Projevi se nekdy i S\_gr?}  Context: separate \(S_e\), \(\Sl\),
residual saturations, storage, relative permeability, and measurement proxies.}
\newcommand{\annCommentSeven}{\textbf{Exact available note:}
\emph{tensor pro permeabiliru by asi mel byt otoceny podle orientace beddingu,
mohlo by byt napsane nekde explicitne v uvodni casti.}  Context: make the
bedding-angle and tensor-rotation gate explicit.}
\newcommand{\annCommentEight}{\textbf{Exact available note:}
\emph{explicitni formule pro \(\eps^{sw}\)? nebo neni treba a staci OGS
implementace? V OGS to je jako zmena delta\_sigma\_sw =
-(lambda\_i * pi\_i * S\_eff\string^{lambda\_i - 1}) * delta\_S\_eff / dt, kde
S\_eff je efektivni saturace, tlak -> vG -> S\_eff; lambda\_i a p\_i jsou
efektivne "materialove parametry", jeden pro kazdy smer, mozno kodovat
anizotropii; dt : timestep; prefix delta\_ : prirustky.}  Context: replace the
unsupported swelling-strain shorthand by the active swelling-stress-rate law.}
\newcommand{\annCommentNine}{\textbf{Exact available note:}
\emph{paragraph "Measured value, units, and OGS tie." a analogicke pozdeji: co
znamena "OGS tie"? jakoze chapu v otazce how does (something) tie to OGS, ale
"OGS tie" mi zni divne.}  Context: use "model link" for the measurement role.}
\newcommand{\annCommentTen}{\textbf{Exact available note:}
\emph{Obecne asi chceme smazat vsechno, co zminuje teplotni zavislost.}
Context: remove unnecessary active-temperature-dependence language while retaining
the checked isothermal/TRM evidence boundary.}
\newcommand{\annCommentEleven}{\textbf{Exact available note:}
\emph{"Which pressure Dirichlet conditions are attached in the checked OGS XML?"
Trochu divna odpoved.}  Context: separate the mathematical boundary inventory from
the checked XML attachment gate.}
\newcommand{\annCommentTwelve}{\textbf{Exact available note:}
\emph{"What mechanical condition is used on the niche boundary?" bez explicitni
podminky -> nulove plosne sily, ale jestli to funguje stejne jako linearni model,
tak okrajova podminka by mela byt na effective stress, tj. brat do uvady i
tlakovou okrajovou podminku. je tohle ok v ramci OGS?}  Context: state the
natural mechanical boundary and the remaining hydraulic pressure coupling.}
\newcommand{\annCommentThirteen}{\textbf{Exact available note:}
\emph{"Which parameter family is released in the first workflow?" nerozumim tomu,
co znamena released.}  Context: define "released" as allowed to vary in candidate
inverse runs.}
\newcommand{\annCommentFourteen}{\textbf{Exact available note:}
\emph{"Is permeability an OGS state output?" smazal bych.}  Context: replace the
state-output wording by the input-field comparison rule.}
\newcommand{\annCommentFifteen}{\textbf{Exact available note:}
\emph{"When does an OGS field become a residual?" smazal bych nebo zadefinoval
pojem "residual".}  Context: define residual as model-data mismatch after an
observation operator.}
\newcommand{\annCommentSixteen}{\textbf{Exact available note:}
\emph{"Can it become an absolute residual?" nerozumim pojmu "absolute residual".}
Context: replace it with calibrated likelihood residual: a measurement mismatch
with an observation operator, uncertainty, and posterior contribution.}
\newcommand{\annCommentSeventeen}{\textbf{Exact user request:}
\emph{then look into the new materials: main.tex, DAMH\_ERT20\_COMPACT\_TECHNICAL\_REPORT.tex,
domain\_kl\_100\_3h\_final\_report.tex and carefully make another chapter in the
document and add there these results on inversion, so its compact and properly
tied to the rest of the text.}  Context: add an inversion-results chapter tied to
the ERT-only observation scope and model gates.}

\ifreviewchanges
  \newcommand{\revadd}[3]{%
    \reviewbubble{#1}\textcolor{blue!70!black}{\uline{#2}}%
    \reviewfootnote{#1}{#3}%
  }
  \newcommand{\revdel}[3]{%
    \reviewbubble{#1}\textcolor{red!70!black}{\sout{#2}}%
    \reviewfootnote{#1}{#3}%
  }
  \newcommand{\revreplace}[4]{%
    \reviewbubble{#1}\textcolor{red!70!black}{\sout{#2}}\ %
    \textcolor{blue!70!black}{\uline{#3}}%
    \reviewfootnote{#1}{#4}%
  }
  \newcommand{\revreplaceonly}[4]{%
    \reviewbubble{#1}\textcolor{red!70!black}{\sout{#2}}\ %
    \textcolor{blue!70!black}{\uline{#3}}%
    \reviewfootnote{#1}{#4}%
  }
  \newcommand{\revaddnofn}[2]{%
    \reviewbubble{#1}\textcolor{blue!70!black}{\uline{#2}}%
  }
  \newcommand{\revreplacenofn}[3]{%
    \reviewbubble{#1}\textcolor{red!70!black}{\sout{#2}}\ %
    \textcolor{blue!70!black}{\uline{#3}}%
  }
  \newcommand{\revnote}[2]{%
    \reviewbubble{#1}\reviewfootnote{#1}{#2}%
  }
\else
  \newcommand{\revadd}[3]{#2}
  \newcommand{\revdel}[3]{}
  \newcommand{\revreplace}[4]{#3}
  \newcommand{\revreplaceonly}[4]{}
  \newcommand{\revaddnofn}[2]{#2}
  \newcommand{\revreplacenofn}[3]{#3}
  \newcommand{\revnote}[2]{}
\fi
